\documentclass[12pt,a4paper]{report}

% ── Packages ──────────────────────────────────────────────────────────────────
\usepackage[a4paper, margin=1in]{geometry}
\usepackage{graphicx}
\usepackage{amsmath}
\usepackage{array}
\usepackage{booktabs}
\usepackage{longtable}
\usepackage{multirow}
\usepackage{xcolor}
\usepackage{hyperref}
\usepackage{setspace}
\usepackage{titlesec}
\usepackage{fancyhdr}
\usepackage{tocloft}
\usepackage{enumitem}
\usepackage{caption}
\usepackage{subcaption}
\usepackage{listings}
\usepackage{float}
\usepackage{parskip}
\usepackage[T1]{fontenc}
\usepackage{pdfpages}

% ── Hyperref Setup ────────────────────────────────────────────────────────────
\hypersetup{
    colorlinks=true,
    linkcolor=black,
    urlcolor=black,
    citecolor=black
}

% ── Spacing ───────────────────────────────────────────────────────────────────
\onehalfspacing

% ── Chapter Title Formatting ──────────────────────────────────────────────────
\titleformat{\chapter}[display]
  {\normalfont\Large\bfseries\centering}
  {\chaptertitlename\ \thechapter}{12pt}{\Large}
\titlespacing*{\chapter}{0pt}{-10pt}{20pt}

% ── Listings (code style) ─────────────────────────────────────────────────────
\lstset{
  basicstyle=\ttfamily\small,
  breaklines=true,
  frame=single,
  backgroundcolor=\color{gray!10}
}

% ═══════════════════════════════════════════════════════════════════════════════
\begin{document}
\pagenumbering{gobble}
% ══════════════════════════════════════════════════════════════
%  TITLE PAGE
% ══════════════════════════════════════════════════════════════
\begin{titlepage}
    \centering

    {\large \textbf{K. K. Wagh Institute of Engineering Education \& Research, Nashik}}\\[0.2cm]
    {\small (Autonomous Institute affiliated to Savitribai Phule Pune University)}\\[1.5cm]

   % College Logo
    \includegraphics[width=4cm]{college_logo.png}\\[1cm]

    {\Huge \textbf{Internship Report}}\\[0.5cm]
    {\large On}\\[0.5cm]

    {\Large \textbf{Excel Intelligence System ---}}\\[0.2cm]
    {\Large \textbf{AI-Powered Analytics Platform}}\\[1.5cm]

    {\large By}\\[0.5cm]
    {\large \textbf{[Your Full Name]}}\\[0.3cm]

    Roll No.: [Your Roll No.]\\
    Class: B.Tech AIDS [Your Division]\\[1cm]

    \textbf{Under Guidance of}\\[0.3cm]

    \textbf{[Industry Mentor Name]}\\
    (Designation, MKCL -- Maharashtra Knowledge Corporation Limited)\\[0.5cm]

    and\\[0.5cm]

    \textbf{[Institute Guide Name]}\\
    (Asst. Professor)\\[1cm]

    Department of \textbf{Artificial Intelligence and Data Science Engineering}\\[0.5cm]

    Academic Year: \textbf{2025-26}

\end{titlepage}

% ══════════════════════════════════════════════════════════════
%  FRONT MATTER
% ══════════════════════════════════════════════════════════════
\pagenumbering{roman}
\setcounter{page}{1}

% ── Certificate ───────────────────────────────────────────────
\chapter*{Certificate}
\addcontentsline{toc}{chapter}{Certificate}

This is certified that Internship Report entitled \textbf{``Excel Intelligence System --- AI-Powered Analytics Platform''} submitted by \textbf{``[Your Full Name]''} of B.Tech Artificial Intelligence and Data Science for the work completed at \textbf{``Maharashtra Knowledge Corporation Limited (MKCL), Nashik''} is approved by us for submission. It is also certified that, to the best of our knowledge, the report represents actual work carried out by the student in the industry.

\vspace{3cm}

\begin{minipage}[t]{0.45\textwidth}
\vspace{0.3cm}
\textbf{[Institute Guide Name]}\\
Internship Institute Guide\\
Assistant Professor\\
Department of Artificial Intelligence and Data Science
\end{minipage}
\hfill
\begin{minipage}[t]{0.45\textwidth}
% \includegraphics[width=3cm]{sign.png}\\
\textbf{[Industry Mentor Name]}\\
Internship Industry Guide\\
Designation: [Designation]\\
Name of Industry: MKCL
\end{minipage}

\vspace{3cm}

\begin{minipage}[t]{0.45\textwidth}
\textbf{Prof. Dr. D. V. Medhane}\\
Head,\\
Department of Artificial Intelligence \\
and Data Science

\vspace{1cm}

\textbf{External Examiner}\\
Name: \underline{\hspace{4cm}}\\
Company Name: \underline{\hspace{4cm}}\\
Designation: \underline{\hspace{4cm}}\\
Signature with date: \underline{}\\
\end{minipage}
\hfill
\begin{minipage}[t]{0.45\textwidth}
\textbf{Prof. Dr. K. N. Nandurkar}\\
Director\\
K. K. Wagh Institute of Engineering Education \& Research, Nashik
\end{minipage}

% ── Certificate of Internship ─────────────────────────────────
\newpage
\chapter*{Certificate of Internship}
\addcontentsline{toc}{chapter}{Certificate of Internship}

\begin{figure}[H]
    \centering
    \fbox{\parbox{0.9\textwidth}{\centering\vspace{4cm}\textbf{[Place Offer Letter / Internship Certificate Image Here]}\\[0.5cm]{\small Replace this placeholder with your offer letter image.}\vspace{4cm}}}
    \caption{Offer Letter -- Maharashtra Knowledge Corporation Limited (MKCL)}
    \label{fig:offer_letter}
\end{figure}

% ── Acknowledgement ───────────────────────────────────────────
\newpage
\chapter*{Acknowledgement}
\addcontentsline{toc}{chapter}{Acknowledgement}

First, I would like to express my sincere gratitude to \textbf{Maharashtra Knowledge Corporation Limited (MKCL)} for providing me the opportunity to work on the \textbf{Excel Intelligence System} project during my internship. This experience has been instrumental in shaping my understanding of enterprise AI application development, modern full-stack engineering, and production-quality software systems.

I am deeply thankful to \textbf{[Industry Mentor Name]}, my industry guide at MKCL, for their continuous guidance, mentorship, and encouragement throughout the internship. Their expertise in AI systems and software architecture helped me navigate complex technical challenges and adopt best practices in system design.

I would also like to thank my institute guide, \textbf{[Institute Guide Name]}, Assistant Professor, Department of Artificial Intelligence and Data Science, for their academic support and valuable feedback throughout the internship period.

I extend my gratitude to \textbf{Prof. Dr. D. V. Medhane}, Head of the Department of Artificial Intelligence and Data Science, for facilitating the internship program and providing an enabling academic environment.

I am grateful to \textbf{Dr. K. N. Nandurkar}, Director, K. K. Wagh Institute of Engineering Education \& Research, Nashik, for his institutional support and encouragement.

I sincerely acknowledge the support of all my colleagues at MKCL, my department faculty members, and my peers who contributed to the successful completion of this internship.

\vspace{2cm}
\begin{flushright}
\textbf{[Your Full Name]}\\
B.Tech Artificial Intelligence and Data Science\\
Roll No.\ [Your Roll No.]
\end{flushright}

% ── Declaration ───────────────────────────────────────────────
\newpage
\chapter*{Declaration}
\addcontentsline{toc}{chapter}{Declaration}

This is to certify that Internship Report submitted by me is an outcome of my independent and original work completed at \textbf{Maharashtra Knowledge Corporation Limited (MKCL), Nashik}. I have duly acknowledged all the sources from which the ideas and extracts have been taken. The Internship Report is free from any plagiarism and has not been submitted elsewhere for publication.

\vspace{2cm}
\begin{flushright}
\textbf{[Your Full Name]}\\
B.Tech Artificial Intelligence and Data Science\\
Roll No.\ [Your Roll No.]
\end{flushright}

\vspace{1cm}
\noindent Date: [DD/MM/2026]\\
Place: Nashik

% ── Abstract ──────────────────────────────────────────────────
\newpage
\chapter*{Abstract}
\addcontentsline{toc}{chapter}{Abstract}

This report presents the design, development, and implementation of the \textbf{Excel Intelligence System} --- a full-stack, AI-powered analytics platform built during an internship at \textbf{Maharashtra Knowledge Corporation Limited (MKCL), Nashik}. The platform allows users to upload Excel and CSV files, connect to external databases, and interact with their data through natural language conversations --- automatically generating SQL queries, interactive dashboards, charts, and narrative reports without requiring any technical expertise.

At the core of the system is a novel \textbf{Tri-Agent AI Architecture} consisting of a Metadata Agent, an Analytical Agent, and a Plot Agent, all orchestrated by a central engine. When a user asks a question in plain English, the system intelligently routes the query to the right agent, generates the appropriate SQL, executes it securely, and presents the results as tables, charts, or KPI cards. The platform also features a complete \textbf{ETL (Extract, Transform, Load) pipeline} for ingesting data from external databases, a \textbf{semantic layer} for encoding business definitions, and a \textbf{public sharing system} for distributing dashboards and reports without requiring login.

The backend was developed using \textbf{Python and FastAPI}, while the frontend was built with \textbf{Vue 3 and TypeScript}. Key technologies include \textbf{LangChain} for AI orchestration, \textbf{PostgreSQL with pgvector} for data storage and semantic caching, \textbf{ECharts and GridStack} for interactive dashboards, and \textbf{Langfuse} for AI observability and cost tracking.

\vspace{0.5cm}
\noindent\textbf{Key Words:} Artificial Intelligence, Natural Language Querying, Dashboard Generation, ETL Pipeline, LangChain, FastAPI, Vue 3, PostgreSQL, Semantic Layer, Multi-Tenant Architecture.

% ── Table of Contents ─────────────────────────────────────────
\newpage
\tableofcontents

\newpage
\listoffigures
\addcontentsline{toc}{chapter}{List of Figures}

\newpage
\listoftables
\addcontentsline{toc}{chapter}{List of Tables}

% ══════════════════════════════════════════════════════════════
%  MAIN MATTER
% ══════════════════════════════════════════════════════════════
\pagenumbering{arabic}
\setcounter{page}{1}

% ──────────────────────────────────────────────────────────────
%  CHAPTER 1 – INTRODUCTION
% ──────────────────────────────────────────────────────────────
\chapter{Introduction}

\section{Introduction}

Organizations today accumulate enormous volumes of data --- spread across Excel files, CSV exports, and operational databases. Yet extracting useful, timely insights from this data remains surprisingly difficult. Business users cannot write SQL, manual analysis is repetitive and error-prone, and traditional BI tools require extensive setup and technical expertise. The result is a persistent gap between \textit{having} data and \textit{understanding} it.

Artificial Intelligence, particularly Large Language Models (LLMs), offers a transformative opportunity: enabling anyone to interact with data by simply asking questions in plain English. Rather than writing complex SQL or configuring dashboard widgets manually, users can say \textit{``Show me monthly revenue trends''} or \textit{``Compare sales across regions''} and receive instant, accurate visual answers.

The \textbf{Excel Intelligence System} was developed during the internship at \textbf{Maharashtra Knowledge Corporation Limited (MKCL)} to make this vision a reality. It is an AI-powered analytics platform that turns scattered files and databases into interactive dashboards, charts, and reports --- all driven by natural language. The system is not a generic AI wrapper; it is \textit{grounded in organizational data}, understanding table structures, business definitions, and naming conventions to produce relevant, accurate results.

The platform supports two major workflows:
\begin{enumerate}[leftmargin=1.5cm]
    \item \textbf{File-based Analytics:} Upload an Excel or CSV file, preview its contents, ask questions in natural language, and generate AI-powered dashboards.
    \item \textbf{Workspace-based Analytics:} Connect external databases, run ETL pipelines to import data, define business semantics (metrics, dimensions, synonyms), and perform advanced querying, reporting, and dashboard creation.
\end{enumerate}

\section{Objectives}

The objectives of this internship project are to:
\begin{itemize}[leftmargin=1.5cm]
    \item Build an AI-powered analytics platform that enables non-technical users to query data using natural language.
    \item Design a Tri-Agent AI system that intelligently routes user questions to specialized agents for metadata, analysis, and visualization.
    \item Develop interactive, drag-and-drop dashboards that can be generated by AI with a single click.
    \item Create an ETL pipeline engine to connect external databases, transform data, and load it into analytical workspaces.
    \item Implement a semantic layer allowing users to define business metrics, dimensions, and synonyms so the AI produces organization-specific results.
    \item Build public sharing capabilities so dashboards and reports can be viewed without login.
    \item Integrate LLM observability tools for monitoring AI cost, accuracy, and performance.
    \item Ensure enterprise-grade security through authentication, role-based access control, and multi-tenant data isolation.
    \item Gain practical, industry-level experience across the full AI application development lifecycle --- from design through deployment.
\end{itemize}

\section{Scope for Learning}

The internship provided comprehensive, hands-on learning across the full spectrum of modern AI-powered application development:

\begin{itemize}[leftmargin=1.5cm]
    \item \textbf{AI and LLM Engineering:} Multi-agent system design, prompt engineering, intent classification, semantic caching, Retrieval-Augmented Generation (RAG), and autonomous error correction.
    \item \textbf{Full-Stack Web Development:} Backend API design with Python and FastAPI; frontend development with Vue 3, TypeScript, and modern UI libraries.
    \item \textbf{Data Engineering:} ETL pipeline development, database connectivity, data transformation, and delta synchronization.
    \item \textbf{Dashboard and Visualization Design:} Interactive chart generation, drag-and-drop layouts, KPI card design, and export to PDF/PNG.
    \item \textbf{Database Design:} Multi-tenant schema architecture, relational modeling, vector embeddings for semantic search.
    \item \textbf{DevOps and Deployment:} Docker containerization, automated startup scripts, and production logging.
    \item \textbf{Software Engineering Practices:} Modular architecture, code reviews, version control, debugging, and testing.
\end{itemize}

\section{Organization of the Internship Report}

\textbf{Chapter 1} introduces the project, its objectives, and scope of learning.

\textbf{Chapter 2} provides an overview of MKCL, the organization where the internship was conducted.

\textbf{Chapter 3} describes the system in detail --- the problem it solves, key features, AI architecture, how each module works, and the technologies used.

\textbf{Chapter 4} presents the results achieved, a comparative analysis, and recommendations for future improvement.

\textbf{Chapter 5} provides the conclusion and future scope.

% ──────────────────────────────────────────────────────────────
%  CHAPTER 2 – OVERVIEW OF COMPANY
% ──────────────────────────────────────────────────────────────
\chapter{Overview of the Company}

\section{Introduction}

\textbf{Maharashtra Knowledge Corporation Limited (MKCL)} is a premier public-private partnership organization established by the Government of Maharashtra for transforming education, governance, and citizen services through information technology. MKCL develops and operates large-scale digital platforms serving millions of students, teachers, and citizens across Maharashtra.

\section{Brief History}

MKCL was incorporated in 2007 under the Companies Act as a Section 25 (not-for-profit) company. It was founded with the vision of creating inclusive, technology-driven knowledge ecosystems for Maharashtra.

\noindent\textbf{Key Highlights:}
\begin{itemize}[leftmargin=1.5cm]
    \item \textbf{Established:} 2007, Government of Maharashtra initiative
    \item \textbf{Headquarters:} Pune, Maharashtra (with offices across the state)
    \item \textbf{Domain:} Education Technology, Digital Governance, AI, and Data Analytics
    \item \textbf{Scale:} Serves millions of users through online examination systems, admission portals, scholarship management, and digital literacy programs
    \item \textbf{Innovation:} Active investment in Artificial Intelligence, Machine Learning, Natural Language Processing, and intelligent automation
    \item \textbf{Mission:} Transform education and governance through knowledge management and information technology
\end{itemize}

\section{Brief Overview}

\subsection{Industry Overview}

MKCL operates at the intersection of Education Technology (EdTech), Governance Technology (GovTech), and Enterprise AI. The organization builds and maintains critical public-service platforms --- from online exam systems handling lakhs of concurrent users to centralized admission portals managing university-level intake across Maharashtra.

The organization's commitment to innovation has led to investments in AI-powered analytics and intelligent data systems, enabling smarter institutional decision-making and operational efficiency. The Excel Intelligence System was developed as part of this innovation initiative.

\begin{table}[H]
\centering
\caption{Technology Areas at MKCL}
\begin{tabular}{>{\bfseries}p{4.5cm} p{8cm}}
\toprule
\textbf{Area} & \textbf{Description} \\
\midrule
Education Platforms & Online examinations, admissions, scholarship management \\
AI and Analytics & Natural language querying, intelligent dashboards, data-driven insights \\
Backend Systems & Scalable APIs, microservices, database management \\
Frontend Applications & Modern web applications with responsive, interactive interfaces \\
Cloud and DevOps & Docker containerization, automated deployment, monitoring \\
Data Engineering & ETL pipelines, data warehousing, semantic data modeling \\
\bottomrule
\end{tabular}
\label{tab:mkcl_areas}
\end{table}

\section{About the Internship}

The internship was carried out at \textbf{MKCL, Nashik} under the guidance of \textbf{[Industry Mentor Name]}. The assigned role was \textbf{Full-Stack AI Developer}, working on the design, development, and implementation of the Excel Intelligence System.

The internship covered the complete software development lifecycle:

\begin{enumerate}[leftmargin=1.5cm]
    \item \textbf{Requirements and Design Phase:} Understanding the business problem, studying user needs, exploring existing tools and their limitations, and defining the system architecture.
    \item \textbf{Core Development Phase:} Building the Tri-Agent AI system, backend APIs, frontend interfaces, ETL pipeline, and workspace management system.
    \item \textbf{Dashboard and Visualization Phase:} Developing AI-powered dashboard generation, drag-and-drop layouts, chart systems, and public sharing features.
    \item \textbf{Integration and Optimization Phase:} End-to-end testing, performance optimization, observability integration, and production-readiness improvements.
    \item \textbf{Documentation Phase:} Writing technical documentation, preparing architecture reports, and creating knowledge transfer materials.
\end{enumerate}

% ──────────────────────────────────────────────────────────────
%  CHAPTER 3 – PROBLEM STUDIED / WORK UNDERTAKEN
% ──────────────────────────────────────────────────────────────
\chapter{Problem Studied / Work Undertaken}

\section{Introduction}

This chapter describes the core problem the project addresses, how the system works, the key features developed, and the technologies used. The focus is on \textit{what} the system does and \textit{why}, rather than low-level code details.

\section{Problem Statement}

Organizations, including educational institutions and government bodies, face a common challenge: \textbf{they have lots of data but difficulty understanding it quickly.} Specifically:

\begin{enumerate}[leftmargin=1.5cm]
    \item \textbf{Scattered Data:} Important data lives in dozens of Excel files, CSV exports, and operational databases --- with no single place to explore or query it.
    \item \textbf{Dependency on Technical Teams:} Business users and decision-makers depend on analysts or developers to write SQL queries, build reports, and create dashboards --- creating delays.
    \item \textbf{Repetitive Manual Work:} The same analysis (monthly trends, regional comparisons, KPI tracking) is repeated manually every time.
    \item \textbf{Non-Reusable Dashboards:} Dashboards built for one dataset cannot be reused for similar data across months or departments.
    \item \textbf{Security Concerns:} Using third-party AI tools (like ChatGPT) for data analysis means uploading sensitive organizational data to external servers.
    \item \textbf{No Business Context:} Generic AI tools do not understand organizational terms --- ``Net Revenue'', ``Active Students'', or ``Region Codes'' have specific meanings that off-the-shelf tools cannot interpret.
\end{enumerate}

\begin{table}[H]
\centering
\caption{Gap Analysis --- Why Existing Tools Fall Short}
\begin{tabular}{>{\bfseries}p{4cm} p{8.5cm}}
\toprule
\textbf{Current Approach} & \textbf{Limitation} \\
\midrule
Raw Excel Analysis & Trend analysis is manual; dashboards are not reusable; no natural language support \\
External AI Tools (ChatGPT) & Sensitive data goes to third-party servers; no grounding in organizational schema \\
Generic BI Tools (Tableau, Power BI) & Require extensive modeling and setup; cannot handle ad-hoc natural language questions \\
Manual SQL Queries & Requires technical expertise; slow turnaround; error-prone \\
Email-based Reports & Static snapshots; cannot drill down; no interactivity \\
\bottomrule
\end{tabular}
\label{tab:gap_analysis}
\end{table}

\section{System Overview}

The \textbf{Excel Intelligence System} is an internal AI-powered analytics platform that solves these problems by providing:

\begin{itemize}[leftmargin=1.5cm]
    \item A way to \textbf{upload data files} and instantly explore them
    \item \textbf{Natural language querying} --- ask questions in plain English instead of writing SQL
    \item \textbf{AI-generated dashboards} with charts, KPIs, and summaries --- created in one click
    \item \textbf{ETL pipelines} to import data from external databases
    \item A \textbf{semantic layer} to encode business definitions so the AI ``speaks your language''
    \item \textbf{Public sharing} so dashboards and reports can be viewed by anyone with a link
    \item \textbf{All data stays internal} --- no information leaves the organization's servers
\end{itemize}

\begin{table}[H]
\centering
\caption{System Overview}
\begin{tabular}{>{\bfseries}p{4cm} p{9cm}}
\toprule
\textbf{Attribute} & \textbf{Description} \\
\midrule
System Name & Excel Intelligence System \\
System Type & AI-powered self-service analytics platform \\
Primary Users & Business users, analysts, managers, decision-makers \\
Key Innovation & Tri-Agent AI Architecture for natural language data querying \\
Data Sources & Excel files, CSV files, PostgreSQL, MySQL, MongoDB \\
Output & Interactive dashboards, charts, tables, KPI cards, narrative reports \\
Architecture & Full-stack: Python/FastAPI backend + Vue 3/TypeScript frontend \\
Deployment & Docker-based; all data stays on internal servers \\
\bottomrule
\end{tabular}
\label{tab:sysoverview}
\end{table}

\section{Key Features}

\subsection{Feature 1: Document Intelligence --- Upload and Explore Data}

Users can upload \textbf{CSV and Excel files} (including multi-sheet workbooks) and instantly explore their contents. The system automatically:

\begin{itemize}[leftmargin=1.5cm]
    \item Detects column types (numbers, dates, text) and sanitizes column names
    \item Generates a \textbf{data preview} showing sample rows and column statistics (min, max, mean, count)
    \item Stores the data in a secure PostgreSQL database for fast querying
    \item Tracks upload progress in \textbf{real-time} via a progress bar using Server-Sent Events (SSE)
    \item Supports \textbf{individual sheet extraction} from multi-sheet Excel workbooks
\end{itemize}

\subsection{Feature 2: Natural Language Querying --- Ask Questions in Plain English}

The most powerful feature of the system: users can \textbf{ask questions about their data in plain English} and receive instant answers. Examples:

\begin{itemize}[leftmargin=1.5cm]
    \item \textit{``What are the top 10 customers by revenue?''}
    \item \textit{``Show monthly sales trend for 2025''}
    \item \textit{``Compare performance across all regions''}
    \item \textit{``How many records have missing values?''}
\end{itemize}

The system automatically generates the correct SQL query, executes it, and returns results as a table, chart, or narrative summary --- all without the user needing to know SQL.

\subsection{Feature 3: Tri-Agent AI System --- The Brain Behind the Platform}

At the heart of the system is a novel \textbf{Tri-Agent Architecture} that breaks down each user query into specialized tasks:

\begin{table}[H]
\centering
\caption{Tri-Agent AI System}
\begin{tabular}{>{\bfseries}p{3.5cm} p{4cm} p{5cm}}
\toprule
\textbf{Agent} & \textbf{What It Does} & \textbf{Example Query} \\
\midrule
Metadata Agent & Answers questions about data structure & ``What columns does this table have?'' \\
Analytical Agent & Generates and runs SQL queries for data analysis & ``Total sales by region'' \\
Plot Agent & Creates chart configurations from results & ``Show a bar chart of monthly revenue'' \\
\bottomrule
\end{tabular}
\label{tab:triagent}
\end{table}

A central \textbf{HybridAgent} orchestrator classifies each incoming question into one of four intents --- metadata, analytical, plot, or general conversation --- and routes it to the appropriate agent. This ensures that every question gets the most accurate response.

\textbf{How a query flows through the system:}
\begin{enumerate}[leftmargin=1.5cm]
    \item User types: \textit{``Show me sales by region as a bar chart''}
    \item The system classifies this as a \textbf{plot} intent
    \item It checks the \textbf{semantic cache} --- has a similar question been asked before?
    \item If not, the Analytical Agent reads the table schema and generates SQL: \texttt{SELECT region, SUM(sales) ... GROUP BY region}
    \item The SQL is validated (destructive operations like DELETE are blocked) and executed
    \item The Plot Agent determines the best chart type (bar chart) and maps columns to axes
    \item The result --- data table + chart + generated SQL --- is returned to the user
    \item The result is cached so similar future questions are answered instantly
\end{enumerate}

\subsection{Feature 4: AI-Powered Dashboard Generation}

Users can generate complete, interactive dashboards with a \textbf{single click}. The AI analyzes the dataset and automatically creates:

\begin{itemize}[leftmargin=1.5cm]
    \item \textbf{KPI Cards:} Key metrics like total revenue, average order value, record count
    \item \textbf{Charts:} Bar charts, line charts, pie/donut charts, scatter plots based on data patterns
    \item \textbf{Summaries:} AI-written narrative summaries explaining key trends and insights
    \item \textbf{Tables:} Structured data views with sorting and filtering
\end{itemize}

The dashboard uses a \textbf{drag-and-drop grid layout} (GridStack), so users can rearrange, resize, add, or remove widgets freely. Dashboards can also be:

\begin{itemize}[leftmargin=1.5cm]
    \item \textbf{Saved and reopened} across sessions
    \item \textbf{Exported as PDF or PNG} for sharing via email or presentations
    \item \textbf{Cloned as templates} and applied to different datasets (e.g., reuse a monthly report layout for each new month's data)
    \item \textbf{Used in compare mode} to view two datasets side-by-side
\end{itemize}

\subsection{Feature 5: Projects and Insight Boards}

Users can organize their data and dashboards using:

\begin{itemize}[leftmargin=1.5cm]
    \item \textbf{Projects:} Group related files into folder-like structures with subprojects. Each project can have its own dashboard.
    \item \textbf{Insight Boards:} Standalone multi-screen presentation boards where users can curate widgets from different data sources --- ideal for management reviews and presentations.
    \item \textbf{Publishing:} Boards and project dashboards can be published and shared via public URLs.
\end{itemize}

\subsection{Feature 6: ETL Pipeline --- Import Data from External Databases}

The system includes a complete \textbf{ETL (Extract, Transform, Load)} pipeline for importing data from external databases into the analytics platform:

\begin{table}[H]
\centering
\caption{ETL Pipeline Stages}
\begin{tabular}{>{\bfseries}p{2.5cm} p{10cm}}
\toprule
\textbf{Stage} & \textbf{What Happens} \\
\midrule
Connect & User provides database credentials (PostgreSQL, MySQL, or MongoDB); the system tests the connection and discovers available tables \\
Preview & Sample rows from the source table are shown for validation before importing \\
Transform & Users can write Python scripts to clean or reshape data --- or use \textbf{AI to auto-generate} the transform script from a plain English description \\
Dry Run & The pipeline runs on a small sample (100 rows) to verify everything works before committing \\
Execute & Full data extraction, transformation, and loading into a secure workspace \\
Sync & Incremental updates --- only new or changed data is imported using delta synchronization \\
\bottomrule
\end{tabular}
\label{tab:etl_stages}
\end{table}

Transform scripts run in a \textbf{sandboxed environment} with time limits (30 seconds for dry-run, 120 seconds for full execution), ensuring that buggy or malicious code cannot affect the main system.

\subsection{Feature 7: Workspaces --- Curated Analytical Environments}

Once data is imported via ETL, it lives in a \textbf{Workspace} --- an isolated, curated environment for analytical work. Each workspace provides:

\begin{itemize}[leftmargin=1.5cm]
    \item \textbf{Data Catalog:} Browse all tables, view columns, edit descriptions, and see statistical profiles
    \item \textbf{Semantic Layer:} Define business-specific metrics, dimensions, and synonyms so the AI understands organizational terminology
    \item \textbf{Natural Language Chat:} Ask questions about workspace data with streaming responses
    \item \textbf{Report Generation:} Create reports using AI-driven queries or custom SQL, with AI-generated narrative summaries
    \item \textbf{Interactive Dashboards:} Generate, customize, and save workspace-level dashboards
    \item \textbf{Sharing:} Share workspace dashboards and reports via public links
\end{itemize}

\subsection{Feature 8: Semantic Layer --- Teaching the AI Your Business Language}

A unique capability of this platform is the \textbf{semantic layer}, which lets users encode organizational knowledge so the AI produces accurate, business-relevant results:

\begin{table}[H]
\centering
\caption{Semantic Layer Components}
\begin{tabular}{>{\bfseries}p{3cm} p{4cm} p{5.5cm}}
\toprule
\textbf{Component} & \textbf{Example} & \textbf{How It Helps the AI} \\
\midrule
Metrics & ``Net Profit'' = Revenue $-$ Costs & AI uses the formula when user asks about ``net profit'' \\
Dimensions & ``Region'', ``Quarter'', ``Category'' & AI knows which columns to group by \\
Synonyms & ``clients'' = ``customers'' & AI understands that both terms refer to the same column \\
Table Descriptions & ``Monthly sales data by region'' & AI understands the purpose of each table \\
Column Descriptions & ``Total order value in INR'' & AI produces more meaningful queries \\
Relationships & Orders.customer\_id $\rightarrow$ Customers.id & AI can join tables correctly \\
\bottomrule
\end{tabular}
\label{tab:semantic}
\end{table}

\subsection{Feature 9: Public Sharing}

Dashboards, workspace reports, and even chat interfaces can be \textbf{shared publicly} via tokenized URLs:

\begin{itemize}[leftmargin=1.5cm]
    \item Recipients can view shared content \textbf{without needing to log in} or create an account
    \item Shared views are \textbf{read-only} --- all editing controls are removed
    \item Each share link uses a unique, secure token
    \item Supports sharing of: file dashboards, project dashboards, workspace dashboards, workspace reports, and workspace chat
\end{itemize}

\subsection{Feature 10: AI Observability with Langfuse}

The platform integrates \textbf{Langfuse}, an open-source LLM observability tool, to monitor and optimize AI performance:

\begin{itemize}[leftmargin=1.5cm]
    \item \textbf{Trace Monitoring:} Every AI query is traced --- showing the exact prompts sent to the LLM and the responses received
    \item \textbf{Cost Tracking:} Token usage (input and output) is tracked to estimate and control API costs
    \item \textbf{Latency Analysis:} Response times are measured for each pipeline stage
    \item \textbf{Seamless Access:} Users access the Langfuse dashboard through a built-in SSO bridge --- no separate login needed
    \item \textbf{Configurable:} Langfuse can be enabled or disabled via environment settings
\end{itemize}

\subsection{Feature 11: Multi-Language Support and Dark Mode}

\begin{itemize}[leftmargin=1.5cm]
    \item The interface supports \textbf{multiple languages} including English, Hindi, Marathi, and Devanagari
    \item \textbf{Dark mode} with smooth toggling provides a comfortable experience in all lighting conditions
    \item The UI uses a professional, modern design system with consistent color palettes and typography
\end{itemize}

\section{System Architecture}

The system follows a modern \textbf{three-tier architecture} with clear separation between the user interface, application logic, and data storage:

\begin{table}[H]
\centering
\caption{System Architecture Layers}
\begin{tabular}{>{\bfseries}p{4cm} p{8.5cm}}
\toprule
\textbf{Layer} & \textbf{Description} \\
\midrule
User Interface & Vue 3 web application with interactive dashboards, chat interfaces, and data explorers \\
AI Engine & Tri-Agent system that understands natural language, generates SQL, and creates visualizations \\
Application Server & FastAPI backend that handles authentication, file management, ETL, and all business logic \\
Caching Layer & Redis and pgvector for semantic caching --- similar questions are answered instantly from cache \\
Database Layer & PostgreSQL for data storage; each workspace gets its own isolated schema for security \\
Observability Layer & Langfuse for monitoring AI performance, costs, and accuracy \\
\bottomrule
\end{tabular}
\label{tab:architecture}
\end{table}

\subsection{Multi-Tenant Data Isolation}

A critical architectural decision is \textbf{schema-level multi-tenancy}: each workspace gets its own dedicated PostgreSQL schema. This means:

\begin{itemize}[leftmargin=1.5cm]
    \item Data from one workspace is \textbf{completely invisible} to other workspaces
    \item AI-generated SQL queries are \textbf{automatically scoped} to the correct workspace
    \item Even if the AI makes an error, it cannot accidentally access data from another workspace
    \item This provides enterprise-grade data isolation without the overhead of separate databases
\end{itemize}

\subsection{Semantic Caching}

The system uses a \textbf{dual-layer caching} strategy to avoid redundant AI calls:

\begin{enumerate}[leftmargin=1.5cm]
    \item \textbf{Redis Cache (fast, in-memory):} Checks for semantically similar queries first --- responses are returned in milliseconds
    \item \textbf{pgvector Cache (persistent):} Falls back to PostgreSQL-based vector search for longer-term cache storage
\end{enumerate}

This means if a user asks \textit{``What were total sales last month?''} and another user later asks \textit{``Show me last month's total sales''}, the second query is answered from cache without calling the AI again.

\section{Technologies Used}

\begin{table}[H]
\centering
\caption{Technology Stack}
\begin{tabular}{>{\bfseries}p{4.5cm} p{8cm}}
\toprule
\textbf{Category} & \textbf{Technologies} \\
\midrule
AI and Language Models & LangChain, OpenAI API, Ollama (local LLMs), Prompt Engineering \\
Backend & Python, FastAPI, Uvicorn, SQLAlchemy \\
Frontend & Vue 3, TypeScript, Vite, Pinia (state management) \\
UI and Styling & Tailwind CSS, shadcn-vue components, Lucide icons \\
Charts and Dashboards & ECharts, Chart.js, GridStack (drag-and-drop layouts) \\
Databases & PostgreSQL (with pgvector extension), Redis, MySQL, MongoDB \\
Data Processing & Pandas, PyArrow, openpyxl (Excel processing) \\
Authentication & JWT tokens, bcrypt password hashing, role-based access control \\
Observability & Langfuse (LLM tracing, cost tracking) \\
Code Editing & CodeMirror 6 (in-browser SQL/Python editor) \\
Deployment & Docker, Docker Compose, automated startup scripts \\
Development Tools & VS Code, Postman (API testing), GitHub, ESLint, Prettier \\
\bottomrule
\end{tabular}
\label{tab:techstack}
\end{table}

\section{User Flows}

The following are the primary ways users interact with the system:

\subsection{Flow 1: File-Based Analytics}

\begin{enumerate}[leftmargin=1.5cm]
    \item Register or log in to the platform
    \item Upload a CSV or Excel file (progress shown in real-time)
    \item Preview data --- see rows, columns, and statistical summaries
    \item Ask questions in natural language via the chat interface
    \item View AI-generated answers as tables, charts, or KPI cards
    \item Generate a full dashboard with one click
    \item Customize the dashboard by dragging, resizing, or adding widgets
    \item Save or export the dashboard as PDF/PNG
\end{enumerate}

\subsection{Flow 2: Database-to-Workspace Analytics}

\begin{enumerate}[leftmargin=1.5cm]
    \item Create a new workspace
    \item Go to the ETL page and connect an external database
    \item Preview source tables and select which ones to import
    \item Write a transform script (or let AI generate one)
    \item Run a dry-run to verify the pipeline on sample data
    \item Execute the full ETL job
    \item Open the workspace catalog --- browse tables, add descriptions, define metrics and synonyms
    \item Chat with the workspace data using natural language
    \item Generate reports and dashboards
    \item Share results via public links
\end{enumerate}

\subsection{Flow 3: Public Sharing}

\begin{enumerate}[leftmargin=1.5cm]
    \item Click ``Share'' on any dashboard or report
    \item Copy the generated public URL
    \item Send the URL to stakeholders
    \item Recipients open the link and view the dashboard --- no login required
\end{enumerate}

\section{Work Undertaken --- Development Activities}

\subsection{Backend Development}

The backend was built using \textbf{Python and FastAPI}, structured as a modular application with clean separation of concerns:

\begin{itemize}[leftmargin=1.5cm]
    \item Designed and implemented \textbf{15 modular router modules} covering authentication, file management, AI querying, dashboards, projects, boards, ETL, workspaces, sharing, and health monitoring.
    \item Built the \textbf{HybridAgent orchestration engine} with intent classification, semantic cache integration, and multi-agent query routing using LangChain.
    \item Developed the \textbf{ETL pipeline engine} with support for PostgreSQL, MySQL, and MongoDB source databases, including sandboxed Python transform execution and delta synchronization.
    \item Implemented \textbf{JWT-based authentication} with role-based access control, ensuring only authorized users can access administrative features.
    \item Created the \textbf{DatabaseManager} module for PostgreSQL connection pooling, dynamic schema creation, and multi-tenant workspace isolation.
    \item Built the \textbf{semantic caching layer} with Redis and pgvector for deduplicating similar AI queries.
    \item Integrated \textbf{Langfuse} for LLM observability with SSO bridge, trace monitoring, and configurable toggle.
    \item Implemented \textbf{production logging} with rotating file handlers for comprehensive debugging and audit trails.
\end{itemize}

\subsection{Frontend Development}

The frontend was built using \textbf{Vue 3 with TypeScript and Vite}, providing a modern, responsive user experience:

\begin{itemize}[leftmargin=1.5cm]
    \item Developed \textbf{20+ application pages} including file explorer, chat interface, smart dashboard, ETL wizard, workspace catalog, workspace chat, reports, and 5 public sharing views.
    \item Built a comprehensive \textbf{dashboard system} with GridStack drag-and-drop layouts, ECharts visualizations, KPI cards, chart builder, template cloning, and PDF/PNG export.
    \item Implemented \textbf{real-time streaming} for upload progress, chat responses, and dashboard generation using Server-Sent Events (SSE) and NDJSON.
    \item Created the \textbf{ETL wizard interface} --- a multi-step guided flow for database connection, table preview, transform editing, dry-run, and full execution.
    \item Built the \textbf{workspace management pages} including data catalog, semantic layer editor, relationship manager, chat, reports, and dashboards.
    \item Developed \textbf{Pinia state management stores} for authentication, documents, boards, and global application state with automatic localStorage persistence.
    \item Implemented \textbf{dark mode}, \textbf{multi-language support} (English, Hindi, Marathi), and a consistent design system using Tailwind CSS and shadcn-vue components.
    \item Created the \textbf{API service layer} with Axios, including JWT interceptors for automatic authentication and 401 redirect handling.
\end{itemize}

\subsection{AI System Development}

\begin{itemize}[leftmargin=1.5cm]
    \item Designed and implemented the \textbf{Tri-Agent Architecture} --- Metadata, Analytical, and Plot agents orchestrated by the HybridAgent engine.
    \item Built the \textbf{intent classification system} that routes queries to the correct agent based on natural language understanding.
    \item Developed \textbf{context engineering} techniques --- injecting table schemas, column statistics, semantic layer definitions, and data profiles into LLM prompts for grounded SQL generation.
    \item Implemented \textbf{SQL validation} to prevent destructive operations (DROP, DELETE, INSERT, UPDATE) from being executed.
    \item Built the \textbf{AI dashboard generator} that analyzes datasets and automatically creates multi-widget dashboards with KPIs, charts, and summaries.
    \item Developed the \textbf{AI transform script generator} for ETL pipelines --- users describe their desired transformation in plain English, and the system generates working Python code.
    \item Implemented \textbf{autonomous error correction} with retry logic for handling LLM rate limits and transient failures.
\end{itemize}

\subsection{Testing and Debugging}

\begin{itemize}[leftmargin=1.5cm]
    \item Performed extensive \textbf{API testing using Postman} across all backend endpoints
    \item Debugged and resolved \textbf{LLM context window overflow} issues by optimizing prompt sizes
    \item Fixed \textbf{connection reset errors} in database connection pooling
    \item Resolved \textbf{cross-tenant data isolation} issues in workspace SQL generation
    \item Tested ETL pipelines across \textbf{multiple database types} (PostgreSQL, MySQL, MongoDB)
    \item Performed \textbf{end-to-end user flow testing} across all major features
\end{itemize}

\section{Conclusion of Work Undertaken}

The development of the Excel Intelligence System involved significant, hands-on engineering work across the full stack --- from AI agent design and backend API development to frontend UI engineering and ETL pipeline construction. The result is a comprehensive, production-ready analytics platform that transforms how non-technical users interact with organizational data.

% ──────────────────────────────────────────────────────────────
%  CHAPTER 4 – RESULTS AND DISCUSSIONS
% ──────────────────────────────────────────────────────────────
\chapter{Results and Discussions}

\section{Introduction}

This chapter presents the outcomes achieved during the internship, a comparative evaluation of the system against traditional approaches, and recommendations for future improvements.

\section{Key Results}

\subsection{AI Querying Results}

\begin{itemize}[leftmargin=1.5cm]
    \item The Tri-Agent system \textbf{successfully classifies} user intent across metadata, analytical, and visualization queries, routing each to the appropriate agent.
    \item Generated SQL queries are \textbf{grounded in actual schema metadata} and semantic layer definitions, producing more accurate results than generic AI tools.
    \item The \textbf{semantic cache} reduces redundant LLM calls --- similar questions asked by different users are answered from cache in milliseconds.
    \item \textbf{SQL validation} prevents any destructive operations, ensuring data safety even if the AI generates incorrect queries.
    \item Context engineering with \texttt{df.describe()} profiles and semantic definitions \textbf{significantly improves} query relevance and accuracy.
\end{itemize}

\subsection{Dashboard and Visualization Results}

\begin{itemize}[leftmargin=1.5cm]
    \item AI-generated dashboards \textbf{automatically identify} the most relevant KPIs, trends, and distributions in a dataset.
    \item The drag-and-drop GridStack layout allows users to \textbf{fully customize} dashboard appearance without coding.
    \item Dashboard templates can be \textbf{reused across datasets}, saving significant time for recurring reports (e.g., monthly reviews).
    \item \textbf{PDF and PNG export} enables easy sharing of dashboard snapshots in presentations and emails.
    \item The compare mode allows \textbf{side-by-side analysis} of two datasets for period-over-period or region-wise comparisons.
\end{itemize}

\subsection{ETL Pipeline Results}

\begin{itemize}[leftmargin=1.5cm]
    \item The ETL engine \textbf{successfully connects to} PostgreSQL, MySQL, and MongoDB databases.
    \item \textbf{AI-generated transform scripts} allow non-technical users to define data transformations in plain English.
    \item \textbf{Sandboxed execution} prevents buggy or malicious transform scripts from affecting the system.
    \item \textbf{Delta synchronization} enables efficient incremental data updates without full re-extraction.
    \item The \textbf{dry-run feature} catches errors early before committing to full data import.
\end{itemize}

\subsection{Workspace and Semantic Layer Results}

\begin{itemize}[leftmargin=1.5cm]
    \item Workspace-level \textbf{schema isolation} provides complete data security between tenants.
    \item The semantic layer \textbf{significantly improves} AI query accuracy by grounding the LLM in organizational terminology.
    \item \textbf{Report sharing} via tokenized URLs enables stakeholders to access analytical results without platform accounts.
    \item \textbf{Auto-profiling} of workspace tables prepares data for AI querying immediately after ETL import.
\end{itemize}

\section{Comparative Analysis}

\begin{table}[H]
\centering
\caption{Comparison: Traditional Approach vs. Excel Intelligence System}
\begin{tabular}{>{\bfseries}p{3.5cm} p{4cm} p{5cm}}
\toprule
\textbf{Aspect} & \textbf{Traditional Approach} & \textbf{Excel Intelligence System} \\
\midrule
Querying Data & Write SQL manually; requires technical expertise & Ask questions in plain English \\
Creating Dashboards & Manual setup in BI tools; hours of configuration & One-click AI generation; drag-and-drop customization \\
Understanding Data & Read documentation; ask a DBA & AI-powered chat explains structure and trends \\
Importing External Data & Custom ETL scripts; manual monitoring & Guided wizard with AI transform generation \\
Reusing Reports & Copy-paste; manual editing & Clone dashboard templates across datasets \\
Sharing Results & Screenshots; email attachments & Live, interactive public URLs \\
Data Security & Data may leave organization & All data stays on internal servers \\
Business Context & Generic; no org-specific understanding & Semantic layer encodes metrics, synonyms, definitions \\
Cost Visibility & No insight into AI usage costs & Langfuse tracks tokens, latency, and costs \\
\bottomrule
\end{tabular}
\label{tab:comparative}
\end{table}

\section{Value Delivered}

\begin{table}[H]
\centering
\caption{Value Delivered by the Excel Intelligence System}
\begin{tabular}{>{\bfseries}p{3.5cm} p{9cm}}
\toprule
\textbf{Value Area} & \textbf{Impact} \\
\midrule
Strategic & Builds internal analytics capability; reduces dependence on external tools and vendors \\
Operational & Saves analyst time; automates repetitive reporting; enables self-service analytics \\
Governance & All data stays internal; role-based access; complete audit trail via Langfuse \\
Knowledge & Business logic is encoded in the semantic layer; dashboards become reusable organizational assets \\
Accessibility & Non-technical users can query data directly; multi-language support broadens reach \\
\bottomrule
\end{tabular}
\label{tab:value}
\end{table}

\section{Professional Development Outcomes}

\begin{itemize}[leftmargin=1.5cm]
    \item \textbf{Full-Stack Proficiency:} Gained comprehensive skills in Python (FastAPI), TypeScript (Vue 3), SQL (PostgreSQL), and AI (LangChain) --- the complete modern AI application stack.
    \item \textbf{System Design Thinking:} Learned to design modular, scalable, multi-tenant architectures with clean separation of concerns.
    \item \textbf{AI Engineering:} Developed practical expertise in prompt engineering, multi-agent orchestration, context engineering, and RAG-style AI systems.
    \item \textbf{Production Mindset:} Adopted production-grade practices including JWT security, structured logging, error handling, and deployment automation.
    \item \textbf{Problem-Solving:} Addressed complex challenges including LLM context window overflow, cross-tenant data isolation, real-time streaming, and ETL sandboxing.
    \item \textbf{Communication:} Improved technical communication through documentation writing, architecture reports, and knowledge transfer sessions.
\end{itemize}

\section{Recommendations}

Based on the internship experience, the following recommendations are offered:

\begin{enumerate}[leftmargin=1.5cm]
    \item \textbf{Automated Testing:} Implement comprehensive test suites for both backend and frontend to catch regressions early, especially in the AI pipeline.
    \item \textbf{Database Migrations:} Adopt a migration framework (like Alembic) for versioned, reversible schema changes instead of ad-hoc table creation.
    \item \textbf{Rate Limiting:} Add API rate limiting to prevent abuse and ensure fair resource allocation.
    \item \textbf{CI/CD Pipeline:} Establish continuous integration and deployment for automated testing, linting, and container building.
    \item \textbf{User Feedback Loop:} Add a feedback mechanism on AI-generated answers so users can flag incorrect results, enabling continuous improvement.
    \item \textbf{Onboarding Guides:} Create step-by-step onboarding tutorials for new users and developers.
\end{enumerate}

% ──────────────────────────────────────────────────────────────
%  CHAPTER 5 – CONCLUSION AND FUTURE SCOPE
% ──────────────────────────────────────────────────────────────
\chapter{Conclusion and Future Scope}

\section{Conclusion}

The internship at Maharashtra Knowledge Corporation Limited (MKCL) provided a rare opportunity to design, develop, and deliver a \textbf{production-grade, AI-powered analytics platform} from the ground up. The \textbf{Excel Intelligence System} directly addresses a real and widespread problem: organizations have abundant data but lack the tools and expertise to understand it quickly.

The \textbf{Tri-Agent AI Architecture} --- with its specialized Metadata, Analytical, and Plot agents --- demonstrated that multi-agent AI systems can deliver significantly more accurate and useful results than monolithic, single-prompt approaches. By grounding the AI in actual schema metadata, statistical profiles, and user-defined semantic layers, the system produces responses that are relevant to the organization's specific context and terminology.

The \textbf{ETL pipeline} with sandboxed execution and AI-generated transform scripts makes data ingestion accessible to non-technical users, while \textbf{workspace-level multi-tenancy} ensures complete data isolation and security. The \textbf{interactive dashboard system} with drag-and-drop layouts, template reuse, and one-click AI generation transforms how users create and consume analytical reports. The \textbf{public sharing} feature extends the system's reach beyond logged-in users, enabling stakeholders to access insights through simple URLs.

Beyond the technical achievements, this internship developed critical professional skills --- system design thinking, AI engineering, full-stack development, production-grade coding practices, debugging complex distributed systems, and technical communication. The experience of building a complete platform spanning backend APIs, frontend interfaces, AI agents, ETL pipelines, and deployment infrastructure provided an unparalleled learning opportunity.

The Excel Intelligence System represents a meaningful step toward making data intelligence accessible to everyone in an organization --- not just the technical few. It is an \textbf{internal intelligence layer over organizational data}, reducing dependence on external tools, manual analysis, and SQL expertise while preserving data security, business context, and reusable analytical workflows.

\section{Future Scope}

The foundation built during this internship opens several directions for further development:

\begin{itemize}[leftmargin=1.5cm]
    \item \textbf{Advanced AI Agents:} Develop multi-step reasoning agents using LangGraph that can plan complex analytical tasks, use tools, and self-correct --- enabling questions like \textit{``Analyze last quarter's performance, identify underperforming regions, and suggest corrective actions.''}

    \item \textbf{Multi-Modal Data Support:} Extend the platform to accept PDFs, Word documents, and images alongside Excel/CSV, using document AI and OCR for extracting structured data from unstructured sources.

    \item \textbf{Real-Time Data Streaming:} Add support for real-time data ingestion via Kafka or Change Data Capture (CDC), enabling live dashboards that update as source databases change.

    \item \textbf{Collaborative Features:} Implement real-time collaboration on dashboards using WebSockets --- shared cursors, live commenting, and version history, similar to Google Docs for analytics.

    \item \textbf{Mobile Application:} Develop a Progressive Web App (PWA) or native mobile app for dashboard consumption on phones and tablets, enabling on-the-go decision-making.

    \item \textbf{Custom LLM Fine-Tuning:} Fine-tune open-source language models (LLaMA, Mistral) on organizational data patterns to improve SQL generation accuracy and reduce dependence on commercial AI APIs.

    \item \textbf{Data Governance:} Add data lineage tracking, data quality monitoring, PII detection, and compliance features for organizations with strict regulatory requirements.

    \item \textbf{Natural Language Dashboard Creation:} Enable users to describe entire dashboard layouts in plain English (e.g., \textit{``Create a dashboard with revenue trends on top, regional comparison in the middle, and top-10 products at the bottom''}).

    \item \textbf{Scheduled Reports:} Add the ability to schedule automated report generation and email delivery --- e.g., ``Send the sales summary dashboard to managers every Monday morning.''

    \item \textbf{Cloud-Native Deployment:} Containerize all services with Kubernetes for auto-scaling, health monitoring, and CI/CD pipelines on Azure or AWS.
\end{itemize}

% ══════════════════════════════════════════════════════════════
%  REFERENCES
% ══════════════════════════════════════════════════════════════
\chapter*{References}
\addcontentsline{toc}{chapter}{References}

\begin{enumerate}[leftmargin=1.5cm, label={[\arabic*]}]
    \item LangChain. (2024). \textit{LangChain Documentation}. \url{https://docs.langchain.com/}
    \item OpenAI. (2024). \textit{OpenAI API Reference}. \url{https://platform.openai.com/docs/}
    \item FastAPI. (2024). \textit{FastAPI Documentation}. \url{https://fastapi.tiangolo.com/}
    \item Vue.js. (2024). \textit{Vue 3 Documentation}. \url{https://vuejs.org/}
    \item PostgreSQL. (2024). \textit{PostgreSQL Documentation}. \url{https://www.postgresql.org/docs/}
    \item pgvector. (2024). \textit{pgvector: Open-Source Vector Similarity Search for PostgreSQL}. \url{https://github.com/pgvector/pgvector}
    \item Langfuse. (2024). \textit{Langfuse: Open-Source LLM Observability}. \url{https://langfuse.com/docs}
    \item Redis. (2024). \textit{Redis Documentation}. \url{https://redis.io/docs/}
    \item ECharts. (2024). \textit{Apache ECharts Documentation}. \url{https://echarts.apache.org/}
    \item GridStack. (2024). \textit{GridStack.js: Dashboard Layout Engine}. \url{https://gridstackjs.com/}
    \item Tailwind CSS. (2024). \textit{Tailwind CSS Documentation}. \url{https://tailwindcss.com/docs}
    \item Docker. (2024). \textit{Docker Documentation}. \url{https://docs.docker.com/}
    \item Pandas. (2024). \textit{Pandas Documentation}. \url{https://pandas.pydata.org/docs/}
    \item Lewis, P., et al. (2020). \textit{Retrieval-Augmented Generation for Knowledge-Intensive NLP Tasks}. Advances in Neural Information Processing Systems (NeurIPS 2020).
    \item Vaswani, A., et al. (2017). \textit{Attention Is All You Need}. Advances in Neural Information Processing Systems (NeurIPS 2017).
    \item Pinia. (2024). \textit{Pinia: The Intuitive Store for Vue.js}. \url{https://pinia.vuejs.org/}
    \item Microsoft. (2024). \textit{TypeScript Documentation}. \url{https://www.typescriptlang.org/docs/}
    \item SQLAlchemy. (2024). \textit{SQLAlchemy Documentation}. \url{https://docs.sqlalchemy.org/}
\end{enumerate}

% ══════════════════════════════════════════════════════════════
%  APPENDIX
% ══════════════════════════════════════════════════════════════
\newpage
\chapter*{Appendix}
\addcontentsline{toc}{chapter}{Appendix}

\section*{A. Feature Summary}

\begin{longtable}{>{\bfseries}p{1cm} p{4cm} p{7.5cm}}
\caption{Complete Feature Summary of the Excel Intelligence System} \\
\toprule
\textbf{No.} & \textbf{Feature} & \textbf{Description} \\
\midrule
\endfirsthead
\toprule
\textbf{No.} & \textbf{Feature} & \textbf{Description} \\
\midrule
\endhead
1 & File Upload and Preview & Upload CSV/Excel files; real-time progress; multi-sheet extraction; data preview with statistics \\
\midrule
2 & Natural Language Querying & Ask questions in plain English; automatic SQL generation, execution, and result formatting \\
\midrule
3 & Tri-Agent AI System & Specialized Metadata, Analytical, and Plot agents with intelligent intent routing \\
\midrule
4 & AI Dashboard Generation & One-click dashboard creation with KPIs, charts, tables, and narrative summaries \\
\midrule
5 & Drag-and-Drop Dashboards & GridStack-powered layout; resize, rearrange, add, and remove widgets freely \\
\midrule
6 & Dashboard Templates & Clone and reuse dashboard layouts across different datasets \\
\midrule
7 & Chart Builder & Custom chart creation with dimension/measure selection and aggregation options \\
\midrule
8 & Dataset Comparison & Side-by-side comparison of two datasets for period-over-period analysis \\
\midrule
9 & Projects and Subprojects & Organize files into folder-like structures with project-level dashboards \\
\midrule
10 & Insight Boards & Multi-screen presentation boards with curated widgets from multiple data sources \\
\midrule
11 & ETL Pipeline & Connect external databases; extract, transform, and load data into workspaces \\
\midrule
12 & AI Transform Generation & Describe desired data transformations in English; AI generates Python code \\
\midrule
13 & Sandboxed Execution & Transform scripts run in isolated environments with time limits for safety \\
\midrule
14 & Delta Synchronization & Incremental data updates via high-water-mark tracking \\
\midrule
15 & Workspaces & Isolated analytical environments with their own tables, metadata, and dashboards \\
\midrule
16 & Semantic Layer & Define business metrics, dimensions, and synonyms to improve AI accuracy \\
\midrule
17 & Workspace Chat & Natural language querying over workspace data with streaming responses \\
\midrule
18 & Report Generation & AI-driven or custom SQL reports with narrative summaries \\
\midrule
19 & Public Sharing & Token-based URLs for sharing dashboards and reports without login \\
\midrule
20 & Authentication and RBAC & JWT-based login with role-based access control \\
\midrule
21 & Multi-Tenant Isolation & Schema-level data isolation between workspaces \\
\midrule
22 & Semantic Caching & Redis + pgvector dual-layer cache for instant responses to similar queries \\
\midrule
23 & Langfuse Observability & LLM trace monitoring, token tracking, cost analysis, and SSO bridge \\
\midrule
24 & Dark Mode & Class-based dark/light theme toggling \\
\midrule
25 & Multi-Language Support & English, Hindi, Marathi, Devanagari \\
\midrule
26 & PDF/PNG Export & Export dashboards as PDF or image files \\
\midrule
27 & Saved Questions & Bookmark frequently asked questions for quick reuse \\
\midrule
28 & Chat History & Persistent conversation history with search and soft-delete \\
\midrule
29 & Auto-Profiling & Automatic statistical profiling of tables after ETL import \\
\midrule
30 & Data Catalog & Browse tables, columns, descriptions, and relationships in workspaces \\
\bottomrule
\end{longtable}

\newpage
\section*{B. Technology Stack Summary}

\begin{table}[H]
\centering
\caption{Technology Stack by Category}
\begin{tabular}{>{\bfseries}p{4.5cm} p{8cm}}
\toprule
\textbf{Category} & \textbf{Technologies Used} \\
\midrule
AI and LLM & LangChain, OpenAI API, Ollama, Prompt Engineering, RAG, Context Engineering \\
Backend Framework & Python, FastAPI, Uvicorn, Pydantic \\
Frontend Framework & Vue 3, TypeScript, Vite \\
State Management & Pinia with localStorage persistence \\
Styling and UI & Tailwind CSS, shadcn-vue, Lucide icons \\
Charts & ECharts, Chart.js \\
Dashboard Layouts & GridStack (drag-and-drop) \\
Primary Database & PostgreSQL with pgvector extension \\
Caching & Redis (in-memory), pgvector (persistent) \\
External DB Support & MySQL, MongoDB, PostgreSQL \\
Data Processing & Pandas, PyArrow, openpyxl \\
Authentication & JWT (PyJWT), bcrypt, passlib \\
Observability & Langfuse \\
Code Editing & CodeMirror 6 \\
Internationalization & vue-i18n (4 languages) \\
Export & jspdf, html2canvas \\
Containerization & Docker, Docker Compose \\
Dev Tools & VS Code, Postman, GitHub, ESLint \\
\bottomrule
\end{tabular}
\label{tab:fullstack_appendix}
\end{table}

\newpage
\section*{C. Learning Outcomes}

The internship provided deep, practical learning in the following areas:

\begin{enumerate}[leftmargin=1.5cm]
    \item \textbf{Multi-Agent AI System Design:} Learned to decompose complex AI tasks into specialized agents, implement intent classification, and orchestrate multi-agent workflows using LangChain.

    \item \textbf{Context Engineering for LLMs:} Developed techniques for enriching AI prompts with schema metadata, statistical profiles, and business definitions to improve SQL generation accuracy.

    \item \textbf{Full-Stack Application Development:} Built a complete application spanning backend APIs (FastAPI), frontend UI (Vue 3), database design (PostgreSQL), and deployment infrastructure (Docker).

    \item \textbf{ETL Pipeline Engineering:} Designed and implemented data pipelines with external database connectivity, sandboxed code execution, and incremental synchronization.

    \item \textbf{Interactive Dashboard Development:} Created drag-and-drop dashboards with real-time AI generation, chart rendering, template reuse, and export capabilities.

    \item \textbf{Semantic Data Modeling:} Implemented a semantic layer allowing business users to encode organizational knowledge (metrics, synonyms, relationships) for AI consumption.

    \item \textbf{Enterprise Security Patterns:} Implemented JWT authentication, RBAC, multi-tenant data isolation, and input validation for production-grade security.

    \item \textbf{Real-Time Web Technologies:} Built streaming interfaces using Server-Sent Events and NDJSON for upload progress, chat responses, and dashboard generation.

    \item \textbf{LLM Observability:} Integrated Langfuse for monitoring AI traces, token usage, latency, and costs --- essential for responsible AI deployment.

    \item \textbf{Production Engineering:} Gained experience with production logging, automated deployment scripts, Docker containerization, and environment configuration management.
\end{enumerate}

% ══════════════════════════════════════════════════════════════
%  PHOTOGRAPH OF MENTOR VISIT
% ══════════════════════════════════════════════════════════════
\newpage
\chapter*{Photograph of Internship Mentor Visit}
\addcontentsline{toc}{chapter}{Photograph of Internship Mentor Visit}

\begin{figure}[H]
    \centering
    % \includegraphics[width=0.80\textwidth, keepaspectratio]{visit.jpeg}
    \fbox{\parbox{0.75\textwidth}{\centering\vspace{5cm}\textbf{[Place Mentor Visit Photograph Here]}\\[0.5cm]{\small Replace this placeholder with the actual photograph.}\vspace{5cm}}}
    \caption{Review Visit -- MKCL, Nashik}
    \label{fig:mentor_visit}
\end{figure}

% ══════════════════════════════════════════════════════════════
\end{document}
